\documentclass[11pt, a4paper]{article}

\usepackage[utf8]{inputenc}
\usepackage[T1]{fontenc}
\usepackage{amsmath, amssymb, amsthm, mathtools}
\usepackage{bm}
% indicator function via boldsymbol (no extra package needed)
\usepackage{booktabs}
\usepackage{graphicx}
\usepackage{hyperref}
\usepackage[margin=2.5cm]{geometry}
\usepackage{enumitem}
\usepackage{xcolor}
\usepackage{tcolorbox}
\usepackage{float}

% Theorem environments
\theoremstyle{definition}
\newtheorem{definition}{Definition}[section]
\newtheorem{example}{Example}[section]
\theoremstyle{plain}
\newtheorem{theorem}{Theorem}[section]
\newtheorem{proposition}{Proposition}[section]
\newtheorem{corollary}{Corollary}[section]
\theoremstyle{remark}
\newtheorem{remark}{Remark}[section]

% Colored boxes
\newtcolorbox{keyresult}[1][]{
    colback=blue!5!white,
    colframe=blue!60!black,
    fonttitle=\bfseries,
    title={Key Insight},
    #1
}
\newtcolorbox{warningbox}[1][]{
    colback=red!5!white,
    colframe=red!60!black,
    fonttitle=\bfseries,
    title={Warning},
    #1
}
\newtcolorbox{refbox}[1][]{
    colback=gray!5!white,
    colframe=gray!50!black,
    fonttitle=\bfseries,
    title={Go Deeper},
    #1
}
\newtcolorbox{keyidea}[1][]{
    colback=green!5!white,
    colframe=green!50!black,
    fonttitle=\bfseries,
    title={Core Idea},
    #1
}
\newtcolorbox{prereq}[1][]{
    colback=orange!5!white,
    colframe=orange!60!black,
    fonttitle=\bfseries,
    title={Read Before Coding},
    #1
}
\newtcolorbox{objective}[1][]{
    colback=blue!5!white,
    colframe=blue!60!black,
    fonttitle=\bfseries,
    title={What You Will Build},
    #1
}

% Custom commands
\newcommand{\VaR}{\operatorname{VaR}}
\newcommand{\ES}{\operatorname{ES}}
\newcommand{\CVaR}{\operatorname{CVaR}}
\newcommand{\EV}{\mathbb{E}}
\newcommand{\PV}{\mathbb{P}}
\newcommand{\RV}{\mathbb{R}}
\newcommand{\ind}{\boldsymbol{1}}

\title{%
    \textbf{Extreme Value Theory for Tail Risk} \\[0.5em]
    \large Modeling the Worst Cases With Mathematical Rigor
}
\author{Andr\'es Gonz\'alez Ortega}
\date{March 2026}

\begin{document}
\maketitle

\begin{abstract}
An applied, example-driven treatment of Extreme Value Theory (EVT) for financial risk.
We start from the practical failure of standard models at extreme quantiles,
develop the two main branches of EVT --- block maxima (GEV) and peaks over threshold (GPD) ---
and derive the EVT-based VaR and ES formulas that allow principled extrapolation into the deep tail.
A worked numerical example runs through the entire document: fitting GPD to S\&P~500 losses
and computing risk measures at the 99.9\% level.
We assume strong probability/statistics background; the focus is on \emph{why} EVT is the right tool for tail risk
and \emph{how} to apply it.
\end{abstract}

\tableofcontents
\newpage

%% ============================================================================
\section{Why We Need EVT: The Extrapolation Problem}
\label{sec:why-evt}
%% ============================================================================

\begin{example}[The 25-sigma day]\label{ex:25sigma}
On October~19, 1987, the S\&P~500 fell 22.6\% in a single session.
The annualized volatility at the time was roughly 20\%, giving a daily standard deviation of about
$\sigma \approx 20\% / \sqrt{252} \approx 1.26\%$.
Under a normal model, a 22.6\% loss is a $22.6 / 1.26 \approx 18$-sigma event
(some estimates reach 25-sigma, depending on the estimation window).

The probability of an 18-sigma event under the normal distribution is approximately $10^{-72}$.
For reference, the universe has existed for about $4.3 \times 10^{17}$ seconds.
Even if we observed one return per \emph{second} since the Big Bang,
an 18-sigma event would still be impossible under normality.

Yet it happened.
\end{example}

The problem is not that we made a small modeling error.
The problem is structural: the normal distribution (and even the Student-$t$)
was designed to describe the \emph{bulk} of the data --- the central 90--95\% of outcomes.
When we use a model calibrated to the center to extrapolate into the 0.1\% tail,
the errors are not small; they are \emph{catastrophic}.

Consider the concrete stakes:
\begin{itemize}
    \item \textbf{Basel III/IV} requires banks to compute $\ES_{0.975}$, which depends on quantiles beyond 97.5\%.
    \item \textbf{Internal risk limits} at many institutions use $\VaR_{0.999}$
          (the once-in-four-years loss).
    \item \textbf{Stress testing} (CCAR, DFAST) asks: what happens in a tail scenario
          we may never have observed?
\end{itemize}

We need a theory that is \emph{designed} for extremes --- not borrowed from the center.
That theory is Extreme Value Theory.

\begin{keyresult}
EVT does not assume a specific distribution for the full data.
Instead, it provides a \textbf{universal limiting result}: regardless of the parent distribution,
extreme values (under mild regularity conditions) converge to one of a small number of known distributions.
This is the tail analogue of the Central Limit Theorem.
\end{keyresult}

%% ============================================================================
\section{Block Maxima and the GEV Distribution}
\label{sec:gev}
%% ============================================================================

\subsection{The Setup}

Let $L_1, L_2, \ldots$ be a sequence of daily losses (i.i.d.\ or stationary with weak dependence).
Partition them into blocks of length $n$ (e.g., $n = 21$ trading days $\approx$ one month).
Let $M_n = \max(L_1, \ldots, L_n)$ be the block maximum.

The fundamental question: as $n$ grows, what is the limiting distribution of $M_n$
(after suitable normalization)?

\subsection{The Fisher--Tippett--Gnedenko Theorem}

\begin{theorem}[Fisher--Tippett--Gnedenko]\label{thm:ftg}
If there exist normalizing constants $a_n > 0$ and $b_n \in \RV$ such that
\[
    \PV\!\left(\frac{M_n - b_n}{a_n} \leq x\right) \xrightarrow{d} H(x)
\]
for some non-degenerate distribution $H$, then $H$ must be one of the three extreme value types:
\begin{enumerate}
    \item \textbf{Fr\'echet} ($\xi > 0$): heavy-tailed. $H(x) = \exp(-x^{-1/\xi})$ for $x > 0$.
    \item \textbf{Gumbel} ($\xi = 0$): light-tailed (exponential decay). $H(x) = \exp(-e^{-x})$.
    \item \textbf{Weibull} ($\xi < 0$): bounded upper tail. $H(x) = \exp(-(-x)^{-1/\xi})$ for $x < 0$.
\end{enumerate}
\end{theorem}

These three types unify into a single parametric family.

\begin{definition}[Generalized Extreme Value Distribution]
The \textbf{GEV} distribution has CDF:
\[
    H_{\xi, \mu, \sigma}(x) =
    \exp\!\left\{-\left[1 + \xi\left(\frac{x - \mu}{\sigma}\right)\right]^{-1/\xi}\right\}
\]
for $1 + \xi(x - \mu)/\sigma > 0$, where $\mu \in \RV$ is the location,
$\sigma > 0$ is the scale, and $\xi \in \RV$ is the \textbf{shape} (tail index).
When $\xi = 0$, the expression is interpreted as the limit $\xi \to 0$, giving the Gumbel form.
\end{definition}

The parameter $\xi$ determines everything about the tail:
\begin{table}[H]
\centering
\begin{tabular}{@{}ccll@{}}
\toprule
$\xi$ & \textbf{Type} & \textbf{Tail behavior} & \textbf{Typical domain} \\
\midrule
$> 0$ & Fr\'echet & Heavy tail (polynomial decay) & Equity returns, operational losses \\
$= 0$ & Gumbel & Light tail (exponential decay) & Normal/lognormal data \\
$< 0$ & Weibull & Finite upper endpoint & Bounded physical quantities \\
\bottomrule
\end{tabular}
\caption{The three extreme value types, unified by the shape parameter $\xi$.}
\end{table}

\begin{keyresult}
For financial return data, we almost always find $\hat\xi > 0$ (Fr\'echet domain).
This confirms mathematically what practitioners know empirically:
financial losses have heavier tails than the normal distribution,
and the GEV framework quantifies exactly \emph{how} heavy.
\end{keyresult}

\subsection{Worked Example: Monthly Maxima of S\&P Losses}

\begin{example}[Fitting GEV to block maxima]\label{ex:gev-fit}
Take daily S\&P~500 losses (negative returns, sign-flipped) from 1990--2024 ($\approx$8{,}800 observations).
Partition into monthly blocks ($n \approx 21$ days each), giving $\approx$420 block maxima.

Fitting GEV by maximum likelihood yields (representative values):
\[
    \hat\xi \approx 0.28, \quad \hat\mu \approx 1.52\%, \quad \hat\sigma \approx 0.72\%
\]

\textbf{Interpretation}:
\begin{itemize}
    \item $\hat\xi = 0.28 > 0$: we are in the Fr\'echet domain. The tail is heavy.
      For a Fr\'echet-type distribution with $\xi = 0.28$, only moments up to order $1/\xi \approx 3.6$ are finite
      (the fourth moment does not exist).
    \item $\hat\mu = 1.52\%$: the ``typical'' monthly maximum loss is about 1.5\%.
    \item $\hat\sigma = 0.72\%$: the spread of monthly maxima around the location.
\end{itemize}

From the fitted GEV, the $T$-year return level (the loss exceeded once every $T$ years, on average) is:
\[
    x_T = \hat\mu - \frac{\hat\sigma}{\hat\xi}\left[1 - \left(-\ln\!\left(1 - \frac{1}{12T}\right)\right)^{-\hat\xi}\right]
\]
For $T = 100$ (the once-in-a-century monthly loss): $x_{100} \approx 12.8\%$.
Under a normal model with the same mean and variance, the 100-year return level would be only $\approx 5.1\%$.
The normal underestimates by a factor of 2.5.
\end{example}

\begin{warningbox}
The block maxima approach is \textbf{wasteful}: from each block of 21 observations,
we use only the single largest.
If we have 8{,}800 daily losses but only 420 block maxima,
we have discarded 95\% of the extreme information.
This motivates the peaks-over-threshold approach.
\end{warningbox}

%% ============================================================================
\section{Peaks Over Threshold and the GPD}
\label{sec:gpd}
%% ============================================================================

\subsection{The Idea}

Instead of taking the maximum within each block,
we use \emph{all} observations that exceed a high threshold $u$.
This retains far more data and leads to more efficient estimation.

\begin{definition}[Excess distribution]
For a loss $L$ with distribution function $F$ and a threshold $u < x_F$ (the right endpoint of $F$),
the \textbf{excess distribution function} is:
\[
    F_u(y) = \PV(L - u \leq y \mid L > u), \quad y \geq 0
\]
This is the conditional distribution of the overshoot $L - u$, given that $L$ exceeds $u$.
\end{definition}

\subsection{The Pickands--Balkema--de Haan Theorem}

\begin{theorem}[Pickands--Balkema--de Haan]\label{thm:pbdh}
If $F$ is in the maximum domain of attraction of the GEV with shape parameter $\xi$,
then there exists a positive function $\beta(u)$ such that
\[
    \sup_{0 \leq y < x_F - u} \left| F_u(y) - G_{\xi, \beta(u)}(y) \right| \to 0 \quad \text{as } u \to x_F
\]
where $G_{\xi, \beta}$ is the \textbf{Generalized Pareto Distribution} (GPD):
\[
    G_{\xi, \beta}(y) =
    \begin{cases}
        1 - \left(1 + \frac{\xi y}{\beta}\right)^{-1/\xi} & \text{if } \xi \neq 0 \\[6pt]
        1 - \exp\!\left(-\frac{y}{\beta}\right) & \text{if } \xi = 0
    \end{cases}
\]
with support $y \geq 0$ when $\xi \geq 0$, and $0 \leq y \leq -\beta/\xi$ when $\xi < 0$.
\end{theorem}

\begin{keyidea}
The Pickands--Balkema--de Haan theorem is the \textbf{tail analogue of the CLT}.
Just as sums converge to the normal regardless of the parent distribution,
exceedances over a high threshold converge to the GPD.
The only free parameters are the shape $\xi$ (inherited from the GEV) and the scale $\beta$.
\end{keyidea}

\subsection{Connection Between GEV and GPD}

The two branches of EVT are not independent.
If block maxima follow a $\text{GEV}(\xi, \mu, \sigma)$, then exceedances above a high threshold
follow a $\text{GPD}(\xi, \beta)$ \textbf{with the same shape parameter} $\xi$.
The scale parameters are related by $\beta = \sigma + \xi(u - \mu)$.

This means: fitting GEV to block maxima and fitting GPD to exceedances
should give consistent estimates of $\xi$.
If they disagree substantially, something is wrong (usually the threshold choice).

\subsection{Worked Example: Losses Above the 95th Percentile}

\begin{example}[Fitting GPD to exceedances]\label{ex:gpd-fit}
Continuing with S\&P~500 daily losses (1990--2024, 8{,}800 observations).
Set the threshold $u$ at the 95th percentile of losses: $u \approx 1.65\%$.
This gives $N_u \approx 440$ exceedances (the 5\% largest losses).

Fitting GPD by maximum likelihood to the excesses $\{L_i - u : L_i > u\}$:
\[
    \hat\xi \approx 0.25, \quad \hat\beta \approx 0.58\%
\]

\textbf{Checks}:
\begin{itemize}
    \item $\hat\xi = 0.25$, consistent with the GEV estimate of $\hat\xi = 0.28$ from Example~\ref{ex:gev-fit}. Good.
    \item The GPD fit uses 440 exceedances vs.\ 420 block maxima, giving tighter confidence intervals.
\end{itemize}

The fitted tail distribution for $x > u$ is:
\[
    \bar{F}(x) = \PV(L > x) = \frac{N_u}{n}\left(1 + \hat\xi \cdot \frac{x - u}{\hat\beta}\right)^{-1/\hat\xi}
\]
where $N_u / n = 440 / 8800 = 0.05$ is the empirical exceedance probability.
\end{example}

%% ============================================================================
\section{Threshold Selection}
\label{sec:threshold}
%% ============================================================================

The GPD result is asymptotic: it holds as $u \to x_F$.
In practice, we need a finite threshold that is ``high enough'' for the GPD approximation to be reasonable,
yet ``low enough'' to retain sufficient exceedances for estimation.
This is the classic bias-variance tradeoff of EVT.

\subsection{Mean Residual Life Plot}

\begin{keyidea}
If $L - u \mid L > u \sim \text{GPD}(\xi, \beta)$, then the \textbf{mean excess function}
\[
    e(u) = \EV[L - u \mid L > u] = \frac{\beta}{1 - \xi}
\]
is \textbf{linear in $u$} for thresholds above the point where the GPD approximation kicks in.
The modified scale parameter satisfies $\beta(u) = \beta_0 + \xi u$,
so $e(u) = (\beta_0 + \xi u) / (1 - \xi)$, which is indeed linear with slope $\xi / (1 - \xi)$.

\medskip

\textbf{Mean Residual Life (MRL) plot}: plot the empirical mean excess
\[
    \hat{e}(u) = \frac{1}{N_u}\sum_{i: L_i > u}(L_i - u)
\]
against $u$, with pointwise confidence bands.
Look for the threshold $u_0$ above which the plot becomes approximately linear.
Choose $u \geq u_0$.
\end{keyidea}

\textbf{What the plot looks like in practice (S\&P losses)}:
for small $u$, the MRL plot curves upward (the bulk distribution is not GPD).
Above approximately the 90th--95th percentile, the plot straightens into a roughly linear trend
with positive slope (confirming $\xi > 0$).
Confidence bands widen rapidly for very high thresholds (few exceedances).

\subsection{Parameter Stability Plot}

A complementary diagnostic:

\begin{enumerate}
    \item For a grid of thresholds $u_1 < u_2 < \cdots < u_K$, fit GPD to exceedances above each $u_k$.
    \item Plot $\hat\xi(u_k)$ vs.\ $u_k$ with confidence intervals.
    \item Also plot the \textbf{modified scale} $\hat\beta^*(u_k) = \hat\beta(u_k) - \hat\xi(u_k) \cdot u_k$ vs.\ $u_k$.
\end{enumerate}

If the GPD model is correct above some threshold $u_0$,
both $\hat\xi$ and $\hat\beta^*$ should be \emph{stable} (approximately constant) for $u \geq u_0$.
Instability signals that the threshold is too low.

\subsection{The Hill Estimator}

For distributions in the Fr\'echet domain ($\xi > 0$), the Hill estimator provides a direct estimate
of the tail index from the $k$ largest order statistics $L_{(1)} \geq L_{(2)} \geq \cdots \geq L_{(k)}$:
\[
    \hat\xi_{\text{Hill}}(k) = \frac{1}{k}\sum_{i=1}^{k} \ln L_{(i)} - \ln L_{(k+1)}
\]

\textbf{Hill plot}: plot $\hat\xi_{\text{Hill}}(k)$ vs.\ $k$.
Look for a region where the estimate is stable.
This is the analogue of the parameter stability plot but specific to heavy-tailed data.

\begin{warningbox}
All threshold selection methods involve visual judgment --- there is no fully automatic procedure
that works reliably in all cases.
The recommended practice is to use \emph{all three} diagnostics (MRL plot, parameter stability, Hill plot)
and check for consistency.
If the three methods point to different thresholds, the problem may have structure
(e.g., mixture of two tail regimes) that simple GPD cannot capture.
\end{warningbox}

%% ============================================================================
\section{EVT-Based Risk Measures}
\label{sec:evt-risk}
%% ============================================================================

With a fitted GPD tail, we can compute risk measures at arbitrarily high confidence levels ---
even at quantiles that have no historical observations.

\subsection{EVT VaR}

From the tail estimator in Example~\ref{ex:gpd-fit},
for a confidence level $\alpha$ exceeding the threshold probability $1 - N_u/n$:

\begin{keyidea}
\textbf{EVT-based Value-at-Risk:}
\[
    \VaR_\alpha = u + \frac{\hat\beta}{\hat\xi}
    \left[\left(\frac{n}{N_u}(1-\alpha)\right)^{-\hat\xi} - 1\right]
\]
This inverts the GPD tail estimator at probability level $\alpha$.
\end{keyidea}

\subsection{EVT Expected Shortfall}

By integrating the tail beyond $\VaR_\alpha$:

\begin{keyidea}
\textbf{EVT-based Expected Shortfall} (valid for $\hat\xi < 1$):
\[
    \ES_\alpha = \frac{\VaR_\alpha}{1 - \hat\xi} + \frac{\hat\beta - \hat\xi \, u}{1 - \hat\xi}
\]
Equivalently:
\[
    \ES_\alpha = \frac{\VaR_\alpha + \hat\beta - \hat\xi \, u}{1 - \hat\xi}
\]
The condition $\hat\xi < 1$ ensures the mean of the GPD exists.
For financial data, $\hat\xi \approx 0.2$--$0.4$, so this is always satisfied.
\end{keyidea}

\subsection{Worked Example: 99.9\% Risk Measures}

\begin{example}[EVT vs.\ normal at 99.9\%]\label{ex:evt-risk}
Using the GPD fit from Example~\ref{ex:gpd-fit}:
$u = 1.65\%$, $\hat\xi = 0.25$, $\hat\beta = 0.58\%$, $N_u = 440$, $n = 8800$.

\medskip
\textbf{Step 1: EVT VaR at $\alpha = 0.999$.}
\[
    \VaR_{0.999} = 1.65 + \frac{0.58}{0.25}
    \left[\left(\frac{8800}{440} \times 0.001\right)^{-0.25} - 1\right]
\]
\[
    = 1.65 + 2.32 \left[(0.02)^{-0.25} - 1\right]
    = 1.65 + 2.32 \left[2.659 - 1\right]
    = 1.65 + 2.32 \times 1.659
\]
\[
    \VaR_{0.999}^{\text{EVT}} \approx 5.50\%
\]

\textbf{Step 2: EVT ES at $\alpha = 0.999$.}
\[
    \ES_{0.999} = \frac{5.50 + 0.58 - 0.25 \times 1.65}{1 - 0.25}
    = \frac{5.50 + 0.58 - 0.4125}{0.75}
    = \frac{5.6675}{0.75} \approx 7.56\%
\]

\textbf{Step 3: Compare to normal.}
The sample mean and standard deviation of daily losses are approximately
$\hat\mu_L \approx 0.04\%$ and $\hat\sigma_L \approx 1.15\%$.
Under normality:
\[
    \VaR_{0.999}^{\text{Normal}} = 0.04 + 1.15 \times \Phi^{-1}(0.999)
    = 0.04 + 1.15 \times 3.090 = 3.59\%
\]
\[
    \ES_{0.999}^{\text{Normal}} = 0.04 + 1.15 \times \frac{\phi(3.090)}{0.001} \approx 4.07\%
\]

\medskip

\begin{center}
\begin{tabular}{@{}lcc@{}}
\toprule
\textbf{Risk measure} & \textbf{Normal model} & \textbf{EVT (GPD)} \\
\midrule
$\VaR_{0.999}$ & 3.59\% & 5.50\% \\
$\ES_{0.999}$  & 4.07\% & 7.56\% \\
\bottomrule
\end{tabular}
\end{center}

\medskip

The normal model underestimates $\VaR_{0.999}$ by 35\% and $\ES_{0.999}$ by 46\%.
For a \$100 million portfolio, this is the difference between reserving \$3.6M and \$5.5M against tail losses.
The normal model is not slightly wrong --- it is \textbf{dangerously wrong} in the deep tail.
\end{example}

\begin{keyresult}
The farther into the tail you go, the worse the normal approximation gets.
At 95\%, the normal and EVT estimates are similar.
At 99\%, the gap opens.
At 99.9\%, the normal can underestimate by 30--50\%.
This is because the normal tail decays as $\exp(-x^2/2)$ (Gaussian),
while the heavy-tailed GPD decays as $x^{-1/\xi}$ (polynomial).
The polynomial wins at large $x$.
\end{keyresult}

%% ============================================================================
\section{GARCH + EVT: The McNeil--Frey Approach}
\label{sec:garch-evt}
%% ============================================================================

The EVT analysis above assumed i.i.d.\ losses.
But financial returns exhibit \emph{volatility clustering}:
periods of high volatility persist for weeks or months (see Project~04).
Applying EVT directly to raw returns conflates two effects:
\begin{enumerate}
    \item The \emph{volatility dynamics} (how risk changes over time).
    \item The \emph{tail behavior} of the innovation distribution (how extreme the shocks are).
\end{enumerate}

McNeil \& Frey (2000) proposed a clean two-step procedure.

\subsection{The Method}

\begin{keyidea}
\textbf{Step 1: GARCH filtering.}
Fit a GARCH(1,1) model (or GJR-GARCH, EGARCH) to the return series:
\[
    r_t = \hat\mu_t + \hat\sigma_t z_t
\]
Extract the standardized residuals $\hat{z}_t = (r_t - \hat\mu_t) / \hat\sigma_t$.
These should be approximately i.i.d.\ (check with Ljung-Box test on $\hat{z}_t$ and $\hat{z}_t^2$).

\medskip
\textbf{Step 2: EVT on residuals.}
Apply the GPD to the tails of $\hat{z}_t$ (after sign-flipping for the loss tail).
This gives tail parameters $(\hat\xi_z, \hat\beta_z)$ for the \emph{innovation} distribution.

\medskip
\textbf{Step 3: Conditional risk measures.}
The one-step-ahead conditional VaR is:
\[
    \VaR_{t+1, \alpha} = \hat\mu_{t+1} + \hat\sigma_{t+1} \cdot \VaR_\alpha^{z}
\]
where $\VaR_\alpha^{z}$ is the EVT-based VaR of the standardized residuals and
$\hat\sigma_{t+1}$ is the GARCH forecast of tomorrow's volatility.
Analogously for ES:
\[
    \ES_{t+1, \alpha} = \hat\mu_{t+1} + \hat\sigma_{t+1} \cdot \ES_\alpha^{z}
\]
\end{keyidea}

\subsection{Why This Works}

The key insight is \textbf{separation of concerns}:
\begin{itemize}
    \item GARCH captures the \emph{time-varying volatility} (the conditional scale).
    \item EVT captures the \emph{tail shape} of the innovations (the unconditional distribution of shocks, once volatility is removed).
\end{itemize}

Because the standardized residuals are (approximately) i.i.d.,
the EVT asymptotic theory applies cleanly.
Applying EVT to raw returns, which are serially dependent through volatility clustering,
violates the i.i.d.\ assumption and can lead to biased tail estimates.

\begin{warningbox}
The quality of the GARCH filter matters.
If the GARCH model is misspecified (e.g., using GARCH(1,1) when the data has strong leverage effects),
the standardized residuals will still exhibit dependence,
and the subsequent EVT step inherits that misspecification.
Always check residual diagnostics before applying EVT.
\end{warningbox}

\subsection{Connection to the Project Pipeline}

In the risk-analyst project:
\begin{itemize}
    \item \textbf{Project~04} (Volatility Modeling) provides the GARCH filter:
          fitting GARCH/GJR-GARCH/EGARCH and extracting $\hat\sigma_t$ and $\hat{z}_t$.
    \item \textbf{Project~05} (this project) applies EVT to the filtered residuals.
    \item \textbf{Project~06} (Copula Dependency) uses the same GARCH-filtered residuals
          as input to copula fitting (after probability integral transform).
\end{itemize}

The GARCH + EVT approach produces \emph{conditional} risk measures that adapt to the current
volatility regime while correctly modeling the extreme tail.
This is the state of the art for univariate tail risk estimation.

\begin{refbox}
\begin{itemize}
    \item McNeil, A.J.\ \& Frey, R.\ (2000). ``Estimation of tail-related risk measures for
          heteroscedastic financial time series: an extreme value approach.''
          \emph{Journal of Empirical Finance}, 7(3--4), 271--300.
          The foundational paper for GARCH + EVT.
    \item McNeil, A.J., Frey, R.\ \& Embrechts, P.\ (2015).
          \emph{Quantitative Risk Management: Concepts, Techniques and Tools}. 2nd ed., Ch.~7.
          Textbook treatment.
\end{itemize}
\end{refbox}

%% ============================================================================
\section{References}
\label{sec:references}
%% ============================================================================

The following texts provide the mathematical foundations and practical guidance for EVT in finance.

\begin{enumerate}[label={[\arabic*]}, leftmargin=2cm]
    \item \textbf{de Haan, L.\ \& Ferreira, A.}\ (2006).
          \emph{Extreme Value Theory: An Introduction}.
          Springer.
          The rigorous mathematical treatment.
          Chapters 1--3 cover the GEV and GPD theory;
          Chapter~6 addresses statistical estimation.

    \item \textbf{McNeil, A.J.\ \& Frey, R.}\ (2000).
          ``Estimation of tail-related risk measures for heteroscedastic financial time series:
          an extreme value approach.''
          \emph{Journal of Empirical Finance}, 7(3--4), 271--300.
          The GARCH + EVT methodology.

    \item \textbf{Embrechts, P., Kl\"uppelberg, C.\ \& Mikosch, T.}\ (1997).
          \emph{Modelling Extremal Events for Insurance and Finance}.
          Springer.
          The classic reference for EVT in risk management;
          strong on both theory and applications.

    \item \textbf{Coles, S.}\ (2001).
          \emph{An Introduction to Statistical Modeling of Extreme Values}.
          Springer.
          Accessible, applied treatment with many worked examples.
          Excellent for threshold selection diagnostics.

    \item \textbf{McNeil, A.J., Frey, R.\ \& Embrechts, P.}\ (2015).
          \emph{Quantitative Risk Management: Concepts, Techniques and Tools}. 2nd ed.
          Princeton University Press.
          Chapters 5 and 7 cover EVT and its integration with GARCH.

    \item \textbf{Danielsson, J.}\ (2011).
          \emph{Financial Risk Forecasting}.
          Wiley.
          Chapter~9 provides a practitioner-oriented EVT guide with R code.
\end{enumerate}

\vfill
\hrule
\smallskip
\noindent\textit{Andr\'es Gonz\'alez Ortega --- Risk Analyst Project} \hfill \textit{March 2026}

\end{document}
